\scanpage{28}
Si on a deux cônes \inserted{convexes} \unsure{fermé|fermés} $P,P^*$ dans
$\Phi(M)$, $\Phi(M^*)$, \deleted{qui sont} tels que

a) $P,P^*$ stables par automorphismes

b) $\langle P,P^*\rangle\geq0$

c) Il existe $\varphi,\varphi^*$ polaires l'un de l'autre,
``admissibles'' [i.e. \deleted{\unsure{$\mathrm{Tr}\,u$|Pour}} $\varphi$ non dég. et
$\mathrm{Tr}\,u^\varphi u\geq0$ pour tt $u\in A$], avec $\varphi\in P$, $\varphi^*\in P^*$

alors $P,P^*$ sont néc. deux cônes admissibles
associés, (et en particulier chacun est
exactement le polaire de l'autre, et $\overset{\circ}{P}^{\deleted{\unsure{*}}+}$
consiste exactement en les $\overset{\vee}{\varphi}$ pr $\varphi\in\overset{\circ}{P}^+$).

\uline{Abstraitement}

Nous avons une Algèbre $A$ $\frac12$ simple de type fini sur $Q$,
\inserted{munie d'une forme $\mathrm{tr}$ telle que $\mathrm{tr}(uv)$ \uline{\unsure{non dégénérée|non définie}}}
un \deleted{\unsure{b|bi}}module \inserted{à droite} $\Phi$ sous $A\otimes A$ (élément $\varphi$,
opération de $A$ écrite $u\varphi v$) \deleted{\unsure{biunivoque|inversible|transportés}} libre
de rang $1$ sur $A\otimes1$ et sur $1\otimes A$
\deleted{(\unsure{donc canoniqt|isom. correspond.} $\unsure{\simeq}$ à $A$)}

le \deleted{bimodule} $Q$-dual $\Phi^*$, \deleted{\unsure{églt et isomorphe}} muni de sa
structure de module sur $A^\circ\otimes A^\circ$
\deleted{à $A$ par conséquent}

Pour un élément $\varphi$ de $\Phi$, il revient au $=$
de dire qu'il est une base \deleted{\unsure{à gauche}
à droite, \unsure{Je lui} \unsure{fais} \unsure{on|se} dit \unsure{qu'il est} non dég.}
% archive sheet number: 27
